% Generated draft: requires mathematical and source review.
\documentclass[11pt]{article}
\usepackage[margin=25mm]{geometry}
\usepackage{fontspec}
\setmainfont{Latin Modern Roman}
\usepackage{amsmath,amssymb,mathtools,unicode-math}
\setmathfont{Latin Modern Math}
\usepackage{enumitem,needspace,xurl,hyperref}
\setlength{\parindent}{0pt}
\setlength{\parskip}{5pt}
\setlength{\emergencystretch}{3em}
\newcommand{\N}{\mathbb{N}}
\newcommand{\Z}{\mathbb{Z}}
\newcommand{\Q}{\mathbb{Q}}
\newcommand{\R}{\mathbb{R}}
\newcommand{\C}{\mathbb{C}}
\newcommand{\F}{\mathbb{F}}
\newcommand{\A}{\mathbb{A}}
\newcommand{\PP}{\mathbb{P}}
\newcommand{\E}{\mathbb{E}}
\newcommand{\eps}{\varepsilon}
\newcommand{\dd}{\,\mathrm{d}}
\begin{document}
\section*{UnsolvedMath ID: 30001858}
\textbf{Source problem code:} catalogue entry
\subsection{Definitions and mathematical statement}
The supplied record poses a broad classification question, rather than a specific conjecture. The following gives conventional meanings to its notation; the scope qualifications below remain unresolved.

Let \(K\) be a field, let \(V\) be a finite-dimensional \(K\)-vector space, and let \(\rho\colon G\to\operatorname{GL}_K(V)\) be a linear action of a finite group \(G\). Write \(V^*=\operatorname{Hom}_K(V,K)\) and \(K[V]=\operatorname{Sym}_K(V^*)\), the polynomial coordinate algebra, including when \(K\) is finite. The induced action is determined by \(g\cdot\lambda=\lambda\circ\rho(g)^{-1}\) for \(\lambda\in V^*\). Define
\[
R=K[V]^G=\{f\in K[V]:g\cdot f=f\text{ for every }g\in G\}.
\]
Here a polynomial invariant ring means an isomorphism \(R\cong K[t_1,\ldots,t_m]\) of \(K\)-algebras for some integer \(m\geq 0\). A global complete-intersection presentation means
\[
R\cong K[x_1,\ldots,x_n]/(f_1,\ldots,f_c),
\]
where \(n,c\geq 0\) and \(f_1,\ldots,f_c\) is a regular sequence: its ideal is proper and each \(f_i\) is a non-zero-divisor modulo its predecessors. For a commutative Noetherian ring \(R\), the Cohen--Macaulay property means that \(\operatorname{depth}R_{\mathfrak p}=\dim R_{\mathfrak p}\) for every prime ideal \(\mathfrak p\). The Gorenstein property means that each \(R_{\mathfrak p}\) has finite injective dimension as a module over itself.

\textbf{Requested investigation.} Determine how properties of the finite linear action \(\rho\) characterize or constrain the polynomial, complete-intersection, Gorenstein, and Cohen--Macaulay properties of \(K[V]^G\), including the hypotheses under which particular implications or equivalences hold.

\textbf{Unresolved scope.} The source does not specify the characteristic or other restrictions on \(K\), whether \(\operatorname{char}K\) divides \(|G|\), whether the action must be faithful, or explicit dimension assumptions on \(V\). Finite dimensionality above is a conventional interpretation requiring confirmation. It also does not specify graded versus ungraded polynomial isomorphisms, global versus local complete-intersection conventions, which properties of the action are intended, or a particular implication to establish. Thus this investigation cannot be extracted as a uniquely quantified conjecture from the supplied record; the definitions above do not resolve these source omissions.
\subsection{Short English statement}
The question asks how a finite group's linear action determines whether its polynomial invariant ring is polynomial, a complete intersection, Gorenstein, or Cohen--Macaulay. The supplied formulation leaves the field hypotheses, action assumptions, and intended classification target unspecified, so this is a statement-only research-question draft requiring mathematical review.
\subsection{Sources}
\begin{itemize}
\item Supplied local UnsolvedMath registry record, permanent ID \(30001858\), title \emph{Invariant Ring Structure and Finite Linear Actions}, exact source code \texttt{catalogue entry}; collection: Oberwolfach Reports Open Problems. No individual proposer is supplied. The recorded proposed year is 2011.
\item Supplied source citation: \emph{Computational Group Theory (2012)}, Oberwolfach Reports, pp.\ 2113--2161, DOI \texttt{10.4171/owr/2011/37}, \url{https://doi.org/10.4171/owr/2011/37}. The background describes these reports as documenting mathematical workshops and their open-problem discussions. The supplied literature-source label is ``OWR report 37,'' with the same link: \url{https://doi.org/10.4171/owr/2011/37}.
\item Literature assessment dated exactly \texttt{2026-08-21}: the supplied curation labels the problem open and reports that no general classification resolving this broad invariant-theory question was found. This is supplied dated curation, not an independently verified literature assessment. The background explicitly requests independent verification of time-sensitive status and cited claims.
\item The registry labels the statement \texttt{corrected\_verified} and asserts that the official report confirms the classification question. These are supplied verification claims; the report has not been independently consulted here. The citation's title year 2012, DOI year component 2011, and proposed year 2011 are retained without resolving their relationship.
\end{itemize}
\end{document}
